% 本文件由 paper/generate-content.js 从游戏数据自动生成，请勿手工编辑。

\subsection{赵敏线：朔漠玫瑰}

\subsubsection{第四拍·灵蛇海上}\label{question:zhaomin:4}
\textbf{情节位置：}\hyperref[plot:zhaomin:4]{灵蛇海上完整剧情}。

\textbf{题干：}妊娠早期的常见反应与晕船确有相似之处。请谈谈：如何从中医和西医角度区分妊娠反应与其他原因引起的恶心呕吐？女性在妊娠早期（前三个月）应注意哪些健康事项？

\textbf{对话开场}
\begin{quote}贫道观赵姑娘之状，既似晕船，又似有孕。二者之别，在于——晕船离舟则缓，而妊娠之呕，晨起为甚，且伴有乳胀、倦怠诸症。请问姑娘，近日可还有其他不适？月事可曾按期而至？\end{quote}
\textbf{后台角色提示词}
\begin{quote}\small 你是一位元末明初的医官，精通中医妇科与西域医学。请用半文半白、典雅通俗的语言，与求医者（可能是赵敏这样的蒙古贵族女子）讨论妊娠早期的辨别、调理与注意事项。需涵盖中医（脉诊、饮食禁忌、安胎原则）和西医知识。每次回答控制在150-300字，适时追问引导，但不要一次性说太多。如果求医者回答有深度，可深入探讨；如果回答简单，给予温和提示后继续。\end{quote}
\textbf{判定依据、推导与科普边界}
\begin{quote}妊娠相关恶心呕吐可发生在一天中的任何时段，常伴月经推迟、乳房胀痛、尿频、疲劳或气味敏感；晕船通常与乘船和运动刺激同步，离开环境后缓解。尿或血hCG可提示妊娠，但妊娠位置、孕周与胚胎情况仍需结合病史和超声评估。中医将其归入妊娠恶阻等范畴，须辨证且不能遮蔽脱水、电解质紊乱等风险。早孕期宜按指南补充叶酸、少量多餐、保证水分与休息，避免酒精、烟草和未经医师确认安全的药物或草药，并尽早建立产科保健；若无法进食饮水、体重明显下降、尿量减少、腹痛或阴道流血，应及时就医。\end{quote}
\noindent\textit{资料依据见本页脚注。}
\footnote{Nelson-Piercy C, Dean C, Shehmar M, et al. The management of nausea and vomiting in pregnancy and hyperemesis gravidarum: Green-top Guideline No. 69[J]. BJOG, 2024, 131(7): e1-e30. DOI：\nolinkurl{10.1111/1471-0528.17739}.}
\footnote{American College of Obstetricians and Gynecologists. Morning sickness: nausea and vomiting of pregnancy[EB/OL]. [2026-06-18]. \href{https://www.acog.org/womens-health/faqs/morning-sickness-nausea-and-vomiting-of-pregnancy}{在线链接}.}
\footnote{Centers for Disease Control and Prevention. Folic acid[EB/OL]. 2025-05-20[2026-06-18]. \href{https://www.cdc.gov/folic-acid/index.html}{在线链接}.}
\footnote{傅山. 傅青主女科: 妊娠恶阻[M/OL]. 维基文库, [2026-06-18]. \href{https://zh.wikisource.org/zh-hans/傅青主女科}{在线链接}.}

\subsubsection{第九拍·冰火之岛}\label{question:zhaomin:9}
\textbf{情节位置：}\hyperref[plot:zhaomin:9]{冰火之岛完整剧情}。

\textbf{题干：}中医强调天人合一，认为环境对人体健康有深刻影响。请谈谈：女性在不同环境（如寒冷地区vs炎热地区、干燥vs潮湿）中应如何调理身体？中医因时、因地、因人的三因制宜原则对现代女性健康有何指导意义？

\textbf{对话开场}
\begin{quote}老朽在此冰火同存之地隐居三十载，深有感触。寒为阴邪，易伤阳气，女子体寒者入寒地，则痛经、关节痛加重；热为阳邪，易伤阴液，阴虚者入热地，则潮热盗汗加剧。故《内经》云春夏养阳，秋冬养阴，此因时制宜也。汝可知因地、因人之义乎？\end{quote}
\textbf{后台角色提示词}
\begin{quote}\small 你是一位隐居冰火岛的老医师，精通中医阴阳五行学说和体质辨证。请用半文半白、略带沧桑的语言，与求教者讨论环境气候对女性体质的影响，以及三因制宜（因时、因地、因人）的养生原则。每次回答150-300字。若对方见解有深度，可深入追问；若较简单，给予鼓励后补充新知。\end{quote}
\textbf{判定依据、推导与科普边界}
\begin{quote}三因制宜是中医养生和治疗的核心原则：（1）因时制宜——顺应四时变化（春生夏长秋收冬藏），女性经期、孕期更需注意季节调护；（2）因地制宜——不同地域的气候、水土对人群体质有不同影响（如北方寒冷干燥、南方湿热），迁徙或旅行时需调整饮食起居；（3）因人制宜——每个人的体质（平和质、阳虚质、阴虚质、气虚质、痰湿质、湿热质、血瘀质、气郁质、特禀质）不同，同一环境下的反应也不同。现代女性因工作、旅行等原因频繁迁徙，更需了解自身体质，灵活调适。基础原则：注意保暖避寒（尤其经期）、饮食因地制宜、保持规律作息。\end{quote}
\noindent\textit{资料依据见本页脚注。}
\footnote{佚名. 黄帝内经·素问·宝命全形论[M/OL]. 中国哲学书电子化计划, [2026-06-18]. \href{https://ctext.org/huangdi-neijing/bao-ming-quan-xing-lun/zhs}{在线链接}.}
\footnote{国家市场监督管理总局, 国家标准化管理委员会. GB/T 46939-2025 中医体质分类与判定[S/OL]. 2025-12-31[2026-06-18]. \href{https://openstd.samr.gov.cn/bzgk/std/newGbInfo?hcno=827D6D266BA52A2983F93638DA871028}{在线链接}.}
\footnote{Parent A S, Damdimopoulou P, Johansson H K L, et al. Endocrine-disrupting chemicals and female reproductive health: a growing concern[J]. Nature Reviews Endocrinology, 2025, 21(10): 593-607. DOI：\nolinkurl{10.1038/s41574-025-01131-x}.}
\footnote{Harvard T.H. Chan School of Public Health. How our environment impacts reproductive health[EB/OL]. 2022-09-21[2026-06-18]. \href{https://hsph.harvard.edu/news/how-our-environment-impacts-reproductive-health/}{在线链接}.}

\subsubsection{第十三拍·蝴蝶谷中}\label{question:zhaomin:13}
\textbf{情节位置：}\hyperref[plot:zhaomin:13]{蝴蝶谷中完整剧情}。

\textbf{题干：}中医望闻问切四诊是诊断的基础，其中望诊包括望面色、望舌苔等。作为女性，你可以通过观察面色和舌苔来初步了解自己的身体状况。请谈谈：你对面色与健康的理解是什么？你是否注意过自己或身边女性的面色变化（如面色萎黄、苍白、潮红、晦暗等），这些可能与什么身体状况有关？

\textbf{对话开场}
\begin{quote}蝴蝶谷中，百花盛开。面色者，传统医家视为气血之华；然而灯色、肤色、睡眠与妆容皆会影响所见，断不可望一眼便定病。汝若愿意，可说说近日面色变化是否持续，又伴有头晕、乏力、潮热或其他不适？\end{quote}
\textbf{后台角色提示词}
\begin{quote}\small 你乃是胡青牛再传弟子，隐居于蝴蝶谷中。请用半文半白、深入浅出的语言，与求医者讨论望诊（特别是面色与舌诊）在女性健康中的应用。每次回答150-300字，象胡青牛一样既严谨又温和，适当时追问以了解对方更多症状。\end{quote}
\textbf{判定依据、推导与科普边界}
\begin{quote}望诊具有直观、无创的特点，但面色和舌象受光线、肤色、饮食、药物、口腔卫生等多种因素影响，不能单独用于自我诊断。苍白伴乏力、气促可能提示贫血等问题，应查血常规；突然发绀、意识异常或呼吸困难属于急症信号。中医会把面色、舌象作为辨证线索之一，仍须与闻、问、切合参；现代诊断则需结合病史、体格检查和必要的实验室/影像检查。游戏鼓励观察身体变化，但不鼓励按舌色或脸色自行开药。\end{quote}
\noindent\textit{资料依据见本页脚注。}
\footnote{中国哲学书电子化计划. 难经·六十一难[EB/OL]. \href{https://ctext.org/nan-jing/shen-sheng-gong-qiao/zh}{在线链接}.}
\footnote{中国哲学书电子化计划. 黄帝内经·素问·金匮真言论[EB/OL]. \href{https://ctext.org/huangdi-neijing/jin-gui-zhen-yan-lun/zhs}{在线链接}.}
\footnote{Matos L C, Machado J P, Monteiro F J, Greten H J. Can traditional Chinese medicine diagnosis be parameterized and standardized? a narrative review[J]. Healthcare (Basel), 2021, 9(2):177. PMID:33562368. PMCID:PMC7914658.}
\footnote{Feng L, Xiao W, Wen C, Deng Q. Objectification of tongue diagnosis in traditional medicine, data analysis, and study application[J]. Journal of Visualized Experiments, 2023, (194):e65140. DOI：\nolinkurl{10.3791/65140}.}

\subsubsection{第十七拍·终南问道}\label{question:zhaomin:17}
\textbf{情节位置：}\hyperref[plot:zhaomin:17]{终南问道完整剧情}。

\textbf{题干：}中医养生学是中华民族的瑰宝。隐士所言养阴血、疏肝气、固肾精是女性养生的三大要义。请结合你自身或身边女性的经验，谈谈你对女性养生的理解。你认为现代女性在快节奏生活中，如何将中医养生智慧融入日常？

\textbf{对话开场}
\begin{quote}无量寿福。贫道在这终南山上修行数十载，见的女子病症，十之八九与不知养有关。汝既来问道，贫道且问：汝以为，女子养生之要，与男子有何不同？平日里可有养护自身的习惯？\end{quote}
\textbf{后台角色提示词}
\begin{quote}\small 你是终南山修行的女冠（女道士），精通养生之道。请用带有道家智慧、半文半白的语言，与求教者讨论女性养生之法。回答150-300字。对方若有现代生活烦恼，给予理解的同时点出养生要义；若对方有独到见解，不吝称赞并深入探讨。\end{quote}
\textbf{判定依据、推导与科普边界}
\begin{quote}把中医养生智慧融入现代生活，宜从低风险、可持续的习惯开始：规律睡眠，均衡饮食，避免过度节食，进行每周规律的有氧与抗阻活动，保留稳定的社会支持，并按年龄和风险接受疫苗、筛查与慢病管理。红枣、阿胶、黑芝麻等食物或药材不能笼统称为补血，贫血须查血常规和铁蛋白；中药、艾灸和穴位刺激也有禁忌，应由合格专业人员评估。持续月经异常、剧烈疼痛、异常出血或明显情绪问题，不应只靠养生拖延就医。\end{quote}
\noindent\textit{资料依据见本页脚注。}
\footnote{中国哲学书电子化计划. 黄帝内经·素问·上古天真论[EB/OL]. \href{https://ctext.org/huangdi-neijing/shang-gu-tian-zhen-lun/zhs}{在线链接}.}
\footnote{National Center for Complementary and Integrative Health. Yoga for health: what the science says[EB/OL]. 2024. \href{https://www.nccih.nih.gov/health/providers/digest/yoga-for-health-science}{在线链接}.}
\footnote{National Center for Complementary and Integrative Health. Sleep disorders and complementary health approaches[EB/OL]. \href{https://www.nccih.nih.gov/health/sleep-disorders-and-complementary-health-approaches}{在线链接}.}
\footnote{Li X, Wang C, et al. Qigong and tai chi on human health: an overview of systematic reviews[J]. American Journal of Chinese Medicine, 2022. PMID:36266755.}

\subsection{周芷若线：汉水白莲}

\subsubsection{第四拍·万安寺誓}\label{question:zhouzhiruo:4}
\textbf{情节位置：}\hyperref[plot:zhouzhiruo:4]{万安寺誓完整剧情}。

\textbf{题干：}周芷若在巨大精神压力下出现的身心反应，与现代医学中的功能性下丘脑性闭经（FHA）有相似之处。请从中医和西医角度，阐述精神压力如何影响女性月经周期，以及应如何调理。

\textbf{对话开场}
\begin{quote}贫尼观施主面色——印堂发青，颧红如妆，此乃肝郁化火之象。施主近日是否月事不调，或经前乳胀、烦躁易怒？《女科经纶》云：妇人以血为本，以气为用。气行则血行，气滞则血瘀。施主心中有何郁结，不妨道来。\end{quote}
\textbf{后台角色提示词}
\begin{quote}\small 你是一位在峨眉山修行、精通妇科的医尼。请用半文半白、带有禅意的语言，与求医者（可能是周芷若这样在巨大精神压力下身心失衡的女性）讨论情志与月经的关系及调理之法。每次回答150-300字，温和引导对方倾诉并给予专业建议。\end{quote}
\textbf{判定依据、推导与科普边界}
\begin{quote}功能性下丘脑性闭经（FHA）是由于精神压力、过度节食/运动等因素抑制下丘脑GnRH脉冲式分泌，导致排卵障碍甚至闭经。中医认为思则气结恐则气下怒则气上，强烈或持续的情志刺激可导致：（1）肝气郁结→气滞血瘀→月经后期、痛经、闭经；（2）心脾两虚→气血生化不足→月经量少、色淡；（3）肾虚→冲任亏损→经水断绝。治疗：西医方面以去除诱因（减轻压力、恢复体重、心理治疗）为主，必要时激素治疗；中医以疏肝理气（逍遥散、柴胡疏肝散）、补肾调冲（左归丸、右归丸）、养心安神（甘麦大枣汤）为主。生活调理：规律作息、柔和运动（八段锦）、穴位按摩（太冲、三阴交）。\end{quote}
\noindent\textit{资料依据见本页脚注。}
\footnote{萧埙. 女科经纶[M/OL]. 清代医学古籍. 维基文库: \href{https://zh.wikisource.org/zh-hans/女科經綸}{在线链接}.}
\footnote{Gordon CM, Ackerman KE, Berga SL, et al. Functional hypothalamic amenorrhea: an endocrine society clinical practice guideline[J]. Journal of Clinical Endocrinology \& Metabolism, 2017, 102(5):1413–1439. DOI：\nolinkurl{10.1210/jc.2017-00131}.}
\footnote{Meczekalski B, et al. Functional hypothalamic amenorrhea: impact on bone and neuropsychiatric outcomes[J]. Frontiers in Endocrinology, 2022, 13:953180. DOI：\nolinkurl{10.3389/fendo.2022.953180}.}
\footnote{Endocrine Society. Hypothalamic amenorrhea guideline resources[EB/OL]. \href{https://www.endocrine.org/clinical-practice-guidelines/hypothalamic-amenorrhea}{在线链接}. (accessed 2026-06-18)}

\subsubsection{第九拍·屠狮之会}\label{question:zhouzhiruo:9}
\textbf{情节位置：}\hyperref[plot:zhouzhiruo:9]{屠狮之会完整剧情}。

\textbf{题干：}失眠是许多女性面临的困扰——尤其是经历重大生活事件后。中医将失眠归为不寐范畴，认为与心、肝、脾、肾等脏腑功能失调有关。请谈谈：你在生活中是否经历过失眠？你认为女性的失眠与情绪（如周芷若的愧疚、愤怒、恐惧）有何关联？

\textbf{对话开场}
\begin{quote}阿弥陀佛。施主深夜来访，必是心中有事。《景岳全书》云：不寐证虽有不一，然惟知邪正二字则尽之矣。或以心事抑郁，或以气血不足，或以痰火上扰。贫僧观施主面色，当属肝郁化火、上扰心神——可有心烦易怒、口苦咽干、噩梦纷纭之状？\end{quote}
\textbf{后台角色提示词}
\begin{quote}\small 你是少林药局的医僧，精通不寐（失眠）的辨证论治。请用慈悲温和、半文半白的语言，与求医者（可能是周芷若这样内心充满矛盾与痛苦的人）讨论失眠的原因和调治之法。每次150-300字，循循善诱。\end{quote}
\textbf{判定依据、推导与科普边界}
\begin{quote}失眠（不寐）的发病率女性约为男性的1.5-2倍，与女性激素波动（月经周期、妊娠、更年期）、更高的焦虑抑郁患病率以及社会角色压力有关。中医辨证分型：（1）肝郁化火型——入睡困难、心烦易怒、噩梦（龙胆泻肝汤加减）；（2）心脾两虚型——多梦易醒、心悸健忘（归脾汤）；（3）阴虚火旺型——潮热盗汗、口干心烦（黄连阿胶汤、天王补心丹）；（4）痰热内扰型——胸闷脘痞、噩梦惊恐（温胆汤）。非药物疗法包括：针灸（神门、三阴交、安眠穴）、耳穴压豆、睡前泡脚（加艾叶）、避免睡前使用电子产品（蓝光抑制褪黑素分泌）、保持作息规律。CBT-I（失眠认知行为疗法）是国际首选的非药物治疗。\end{quote}
\noindent\textit{资料依据见本页脚注。}
\footnote{张介宾. 景岳全书［M］. 明代刻本. 电子整理本: 中国哲学书电子化计划（CTEXT）［EB/OL］. \href{https://zh.wikisource.org/zh-hans/景岳全書_(四庫全書本)/卷18}{在线链接}.}
\footnote{Edinger J D, Arnedt J T, Bertisch S M, et al. Behavioral and psychological treatments for chronic insomnia disorder in adults: an American Academy of Sleep Medicine clinical practice guideline［J］. Journal of Clinical Sleep Medicine, 2021, 17(2): 255-262. DOI：\nolinkurl{10.5664/jcsm.8986}.}

\subsubsection{第十五拍·襄阳城墙}\label{question:zhouzhiruo:15}
\textbf{情节位置：}\hyperref[plot:zhouzhiruo:15]{襄阳城墙完整剧情}。

\textbf{题干：}乳腺癌是女性最常见的恶性肿瘤之一。中医认为乳腺疾病与肝气郁结密切相关。请谈谈：你了解哪些乳腺癌的危险因素和早期发现方法？女性应如何进行乳房自我觉察并按风险定期筛查？

\textbf{对话开场}
\begin{quote}大妹子莫要惊慌。汝可知——十之八九的乳房肿块，皆为良性之乳癖（乳腺增生），而非乳岩（癌）。然则既来求医，贫医便细细道来——汝且说说，这肿块何时发现的？是否随月事而有变化（经前胀大、经后缩小）？平日心情如何，可有长期郁结之事？\end{quote}
\textbf{后台角色提示词}
\begin{quote}\small 你是襄阳城中有名的女医，精通乳腺疾病的诊治。请用温和关爱、半文半白的语言，与求医的妇人讨论乳房健康的自我保健知识。语气需缓解其焦虑，同时给予专业的指导。每次150-300字。\end{quote}
\textbf{判定依据、推导与科普边界}
\begin{quote}乳腺癌是女性最常见的恶性肿瘤，WHO数据显示2022年全球约230万女性确诊、约67万人死亡。危险因素：（1）不可改变——女性性别、年龄增长、BRCA1/2等遗传易感、一级亲属家族史、月经初潮早/绝经晚；（2）可改变——饮酒、绝经后肥胖、缺乏运动、长期使用雌孕激素联合HRT等。筛查建议需结合所在国家指南和个人风险：一般风险女性可与医生讨论从40岁起定期乳腺X线摄影（美国USPSTF建议40-74岁每2年一次）；中国指南常结合乳腺X线和超声。高风险女性（如BRCA突变、强家族史）需更早且更密集的MRI/钼靶筛查。日常更推荐乳房自我觉察：熟悉自己的乳房外观和触感，若出现新肿块、乳头凹陷或血性溢液、皮肤橘皮样改变、腋窝肿块等，应及时就医；自我触诊不能替代规范筛查。\end{quote}
\noindent\textit{资料依据见本页脚注。}
\footnote{World Health Organization. Breast cancer fact sheet［EB/OL］. \href{https://www.who.int/news-room/fact-sheets/detail/breast-cancer}{在线链接}.}
\footnote{U.S. Preventive Services Task Force. Breast cancer: screening［EB/OL］. 2024. \href{https://www.uspreventiveservicestaskforce.org/uspstf/recommendation/breast-cancer-screening1}{在线链接}.}
\footnote{国家卫生健康委员会. 乳腺癌诊疗指南（2022年版）［EB/OL］. \href{https://www.nhc.gov.cn/cms-search/downFiles/c001a73dfefc4ace889a1ea6e0230865.pdf}{在线链接}.}
\footnote{Sung H, Ferlay J, Siegel R L, et al. Global cancer statistics 2020: GLOBOCAN estimates of incidence and mortality worldwide for 36 cancers in 185 countries［J］. CA: A Cancer Journal for Clinicians, 2021, 71: 209-249. DOI：\nolinkurl{10.3322/caac.21660}.}
